\section{Worked example: exact rational bisection for \texorpdfstring{$y^3+y=x$}{y3+y=x}}
\label{sec:cubic}

A framework report should include at least one complete cross-layer example.  We use the monotone cubic
\[
 f(y)=y^3+y.
\]
For rational input \(x\) and requested step count \(N\), the goal is to produce an exact rational bracket
\([\ell_N,h_N]\) satisfying
\[
 f(\ell_N)\le x\le f(h_N),
 \qquad
 h_N-\ell_N=2^{-N}(h_0-\ell_0).
\]
This is a finite theorem requiring no completed real root.  Its ordinary interpretation yields a canonical
Cauchy name for the unique real solution, but the iPRA theorem is the displayed rational contract.

\subsection{Coding rational arithmetic}

Represent a rational by an integer numerator and positive denominator.  Every comparison is decided by
cross multiplication:
\[
 \frac ab\le\frac cd
 \quad\Longleftrightarrow\quad
 ad\le cb
 \qquad(b,d>0).
\]
Addition, multiplication, negation, midpoint, and cubic evaluation are primitive recursive on signed
integer codes.  Normalization by gcd is optional for correctness; unreduced fractions keep the program
simpler.  The verifier carries denominator-positivity equations and uses exact integer inequalities.

\subsection{Initial bracket}

Set
\[
 B=|x|+1,
 \qquad \ell_0=-B,
 \qquad h_0=B.
\]
Then
\[
 f(B)=B^3+B\ge B>|x|\ge x,
\]
and
\[
 f(-B)=-B^3-B\le -B<-|x|\le x.
\]
Thus \([\ell_0,h_0]\) brackets the equation.  These are elementary integer inequalities formalizable with
quantifier-free induction for multiplication and exponentiation, or by a small library of ordered-ring
lemmas already instantiated for coded integers.

\subsection{One bisection step}

Given a valid bracket \((\ell,h)\), define the midpoint
\[
 m=\frac{\ell+h}{2}.
\]
The comparison \(f(m)\le x\) is decidable by rational cross multiplication.  Set
\[
 (\ell',h')=
 \begin{cases}
   (m,h),& f(m)\le x,\\
   (\ell,m),& f(m)>x.
 \end{cases}
\]

\begin{lemma}[Bisection preservation]
If \(\ell\le h\) and \(f(\ell)\le x\le f(h)\), then the updated pair satisfies the same bracket
inequalities and
\[
 h'-\ell'=\frac{h-\ell}{2}.
\]
\end{lemma}
\begin{proof}
In the first branch, \(f(m)\le x\le f(h)\), so \((m,h)\) is a bracket.  In the second branch, the decidable
comparison gives \(x<f(m)\), while \(f(\ell)\le x\), so \((\ell,m)\) is a bracket.  In both cases midpoint
arithmetic gives the width identity.  No monotonicity theorem is needed merely to preserve a sign bracket;
monotonicity is used later for uniqueness and error interpretation.
\end{proof}

\subsection{Primitive recursion and invariant}

Pack the rational pair as a state.  Define
\[
  s(0,x)=\langle\ell_0,h_0\rangle,
  \qquad
  s(Si,x)=\operatorname{bisect}(x,s(i,x)).
\]
Use the quantifier-free invariant
\[
 I(i):
 \operatorname{ValidRat}(\ell_i)\wedge
 \operatorname{ValidRat}(h_i)\wedge
 \ell_i\le h_i\wedge
 f(\ell_i)\le x\wedge x\le f(h_i)
 \wedge 2^i(h_i-\ell_i)=h_0-\ell_0.
\]
All rational equalities and inequalities are Booleanized integer equations.  The base proof is the initial
bracket calculation.  The step proof is bisection preservation plus
\[
  2^{i+1}(h_{i+1}-\ell_{i+1})
  =2^{i+1}\frac{h_i-\ell_i}{2}
  =2^i(h_i-\ell_i).
\]
Quantifier-free induction gives \(I(N)\) for arbitrary input step count \(N\).

\begin{theorem}[Finite cubic bisection]
There are fixed primitive-recursive rational functions \(\ell(x,N),h(x,N)\) such that iPRA proves, for
all rational \(x\) and natural \(N\),
\[
 f(\ell(x,N))\le x\le f(h(x,N)).
\]
after decoding the rational inequalities, and
\[
 2^N\bigl(h(x,N)-\ell(x,N)\bigr)=2(|x|+1).
\]
\end{theorem}
\begin{proof}
Compile the trace and invariant above using Theorem~\ref{thm:trace-reflection}.
\end{proof}

\subsection{Uniqueness and represented-root interpretation}

For rationals \(u<v\),
\[
 f(v)-f(u)
 =(v-u)(v^2+uv+u^2+1)>0.
\]
The quadratic factor is positive because
\[
 v^2+uv+u^2
 =\left(v+\frac u2\right)^2+\frac{3u^2}{4}\ge0.
\]
After clearing denominators, this is a finite ordered-ring identity.  Hence \(f\) is strictly increasing
on rational codes, and any two exact roots are equal.  More quantitatively,
\[
 f(v)-f(u)=(v-u)(v^2+uv+u^2+1)\ge v-u
\]
when \(u<v\).  Thus the inverse has the elementary modulus
\[
 |u-v|\le |f(u)-f(v)|
\]
for ordered rationals in the appropriate direction.

Choose the midpoint name
\[
 a(q)=\frac{\ell(x,q+c)+h(x,q+c)}2
\]
with a fixed offset \(c\) large enough to absorb the initial width encoding.  The finite width identity
proves a Cauchy modulus.  Any represented real satisfying \(y^3+y=x\) lies in every bracket and is within
half the bracket width of the midpoint.  The ideal root is therefore a shorthand for the finite bracket
name, not an additional object postulated in iPRA.

\subsection{The staged graph}

The example decomposes into reusable nodes:
\begin{enumerate}
  \item decode and validate the rational input;
  \item compute \(B=|x|+1\) and prove the initial bracket;
  \item define exact midpoint and cubic comparison;
  \item compile the bisection loop from its quantifier-free invariant;
  \item extract the final bracket and width theorem;
  \item optionally prove monotonicity and construct the represented midpoint name.
\end{enumerate}
Each node has a small local arity.  The global theorem may use only the input rational code and \(N\).
Arity-normalized specialization maps each local variable to its concrete global term and composes the
proofs by the static graph theorem.

\subsection{Executable evidence}

The companion file \texttt{code/staged\_contracts.py} implements the exact rational algorithm using
Python's \texttt{Fraction}.  The tests check bracket validity and exact width halving after every step.  The
demo runs ten steps for \(x=7/3\) and writes a replayable JSON record.  The final bracket is
\[
 \frac{275}{256}\le y\le\frac{415}{384},
 \qquad h-\ell=\frac{5}{768}.
\]
These computations are implementation evidence, not the proof of the universal theorem.  Their role is to
exercise the same finite state transition and expose mistakes before the Lean object proof is completed.
